\documentclass[11pt, a4paper]{article}

% Codificación y idioma
\usepackage[T1]{fontenc}
\usepackage[utf8]{inputenc}
\usepackage[spanish]{babel}
\usepackage{csquotes}

% Bibliografía
\usepackage[backend=biber, style=apa]{biblatex}
\addbibresource{referencias_informe.bib}

% Matemáticas
\usepackage{amsmath}
\usepackage{amssymb}

% Tablas y figuras
\usepackage{booktabs}
\usepackage{longtable}
\usepackage{array}
\usepackage{multirow}
\usepackage{graphicx}
\usepackage{float}
\usepackage{caption}
\usepackage{subcaption}
\usepackage{placeins}
\usepackage{threeparttable}
\usepackage{colortbl} % Para \rowcolor

% Listas y enumeraciones
\usepackage{enumitem}

% Formato de página
\usepackage[margin=2.5cm]{geometry}

% Enlaces
\usepackage[colorlinks=true, linkcolor=blue, urlcolor=blue, citecolor=blue]{hyperref}

% Código y diagramas
\usepackage{listings}
\usepackage{tikz}
\usetikzlibrary{shapes.geometric, arrows.meta, positioning, fit}

% Colores
\usepackage{xcolor}
\definecolor{phase0}{HTML}{4CAF50}
\definecolor{phase1}{HTML}{2196F3}
\definecolor{phase2}{HTML}{FF9800}
\definecolor{phase3}{HTML}{9C27B0}

% Encabezados
\usepackage{fancyhdr}
\pagestyle{fancy}
\fancyhf{}
\fancyhead[L]{\small Propuesta de Trabajo — VLM para Dermatología}
\fancyfoot[C]{\thepage}

\title{\textbf{Propuesta de Trabajo: Modelo de Visión--Lenguaje para Diagnóstico Dermatológico y Generación de Reportes Clínicos}}
\author{Facundo Tomás Gaviola \\ \small Universidad Nacional de Cuyo \\ \small Tesina Final de Grado — 2025}
\date{\today}

\begin{document}

\maketitle

\begin{abstract}
Este documento presenta la propuesta de trabajo para el desarrollo y evaluación de un Modelo de Visión--Lenguaje (VLM, Vision-Language Model) orientado al diagnóstico de lesiones cutáneas y la generación automática de reportes clínicos estructurados.
Se detalla el contexto del problema, el marco teórico tecnológico, la justificación de las decisiones de diseño, incluyendo la selección de modelo, datasets, técnicas de ajuste fino y métricas de evaluación y el plan de ejecución en cuatro fases. Todas
las decisiones están respaldadas por un análisis exhaustivo del estado del arte en inteligencia artificial aplicada a dermatología.
\end{abstract}

\tableofcontents
\newpage

% ============================================================================
% I. INTRODUCCIÓN Y PROPÓSITO
% ============================================================================
\section{Introducción y Propósito}

\subsection{Descripción General}

El presente trabajo propone el desarrollo de un sistema basado en un Modelo de Visión--Lenguaje (VLM, Vision-Language Model) capaz de realizar dos tareas complementarias a partir de imágenes dermatológicas: (1) clasificación diagnóstica de lesiones cutáneas,
con énfasis en la detección de melanoma maligno, y (2) generación automática de reportes clínicos estructurados que integren hallazgos visuales, impresión diagnóstica y diagnóstico diferencial.

El modelo seleccionado, LLaVA-Med v1.5 (7B parámetros) \cite{li_llava-med_2023}, será adaptado al dominio dermatológico mediante técnicas de ajuste fino eficiente en parámetros (PEFT, Parameter-Efficient Fine-Tuning) utilizando Adaptación de Bajo Rango 
(LoRA, Low-Rank Adaptation) \cite{hu_lora_2021}. La evaluación se realizará en dos fases: una comparativa frente a modelos de Aprendizaje Profundo (Deep Learning) tradicionales (CNN y Vision Transformer) para clasificación, y una fase de generación de reportes
evaluada con métricas léxicas, semánticas y clínicas.

\subsection{Objetivos}

\subsubsection{Objetivo General}
Desarrollar y evaluar un modelo de visión--lenguaje capaz de asistir en el diagnóstico de lesiones cutáneas y en la generación automática de reportes clínicos estructurados a partir de imágenes dermatológicas, con el fin de mejorar la precisión diagnóstica y
optimizar el flujo de trabajo en el estudio de lesiones de piel.

\subsubsection{Objetivos Específicos}
\begin{enumerate}
    \item Diseñar, implementar y evaluar modelos de detección, incluyendo modelos de Aprendizaje Profundo (Deep Learning) y VLM para identificar posibles melanomas malignos.
    \item Diseñar y entrenar un modelo visión--lenguaje capaz de integrar información visual y textual para producir diagnósticos preliminares y reportes automáticos de lesiones cutáneas.
    \item Evaluar el desempeño del modelo en diferentes escenarios clínicos, analizando su precisión diagnóstica, coherencia textual y utilidad práctica para el flujo de trabajo en dermatología.
\end{enumerate}

% ============================================================================
% II. MARCO DE LA PROBLEMÁTICA
% ============================================================================
\section{Marco de la Problemática}

\subsection{Contexto Global}

Las neoplasias cutáneas representan una carga significativa para la salud pública a nivel mundial. La incidencia del melanoma ha aumentado de manera sostenida durante las últimas tres décadas \cite{khan_ensemble_2025}, constituyendo la forma más agresiva de cáncer de piel. Según datos de la Sociedad Americana del Cáncer, el melanoma localizado presenta una tasa de supervivencia a cinco años del 98\%, cifra que desciende a menos del 20\% cuando se produce metástasis a distancia \cite{khan_ensemble_2025}. Esta brecha subraya la importancia crítica de la detección temprana.

El paradigma diagnóstico actual se basa en la inspección visual y la dermatoscopia, métodos que exhiben una variabilidad interobservador considerable. Las tasas de acuerdo diagnóstico entre dermatólogos experimentados oscilan entre el 65\% y el 80\% dependiendo de la complejidad de la lesión \cite{khan_ensemble_2025}. Esta incertidumbre diagnóstica se ve agravada por disparidades geográficas en el acceso a especialistas, donde las poblaciones rurales enfrentan tiempos de espera que pueden superar los seis meses para una consulta dermatológica.

% IMAGEN SUGERIDA: Ejemplos representativos de las 7 categorías diagnósticas del dataset HAM10000 (MEL, NV, BCC, AKIEC, BKL, DF, VASC).
% Fuente: Tschandl et al. (2018), Fig. 1 — "The HAM10000 dataset, a large collection of multi-source dermatoscopic images of common pigmented skin lesions", Scientific Data.
% Disponible en: https://dataverse.harvard.edu/dataset.xhtml?persistentId=doi:10.7910/DVN/DBW86T

\subsection{Análisis Multidimensional}

Las limitaciones del paradigma diagnóstico actual pueden analizarse desde tres dimensiones interrelacionadas que justifican el desarrollo de herramientas de apoyo basadas en inteligencia artificial.

\textbf{Dimensión médica.} Los subtipos agresivos como el melanoma nodular pueden progresar de estadios tempranos a avanzados en semanas \cite{khan_ensemble_2025}. La detección precoz depende críticamente de la experiencia del profesional y del acceso a herramientas de apoyo diagnóstico. 
Trabajos como el de Esteva et al. (2017) demostraron que las redes neuronales profundas pueden alcanzar un rendimiento diagnóstico comparable al de dermatólogos expertos en tareas de clasificación binaria \cite{esteva_dermatologist_2017}, abriendo la puerta a sistemas de asistencia clínica basados en inteligencia artificial.

\textbf{Dimensión tecnológica.} La evolución desde modelos puramente visuales (CNNs, Vision Transformers) hacia modelos multimodales (VLMs) representa un salto cualitativo: los VLMs no solo clasifican, sino que pueden generar explicaciones textuales de sus decisiones, aumentando la interpretabilidad y la utilidad clínica \cite{hartsock_vision-language_2024}. Sin embargo, el entrenamiento de estos modelos requiere grandes volúmenes de datos pareados (imagen--texto) con validación clínica, un recurso escaso en dermatología.

\textbf{Dimensión social.} Estudios recientes han evidenciado disparidades significativas en el rendimiento de modelos de IA dermatológica según el tono de piel. Daneshjou et al. (2022) demostraron que modelos entrenados predominantemente con imágenes de pieles claras presentan una degradación notable de rendimiento en pieles oscuras (fototipos Fitzpatrick V--VI) \cite{daneshjou_disparities_2022}. El dataset Diverse Dermatology Images (DDI) reveló una diferencia de AUC-ROC de hasta 0.11 entre fototipos claros y oscuros. Esta dimensión refuerza la necesidad de utilizar datasets diversos como Fitzpatrick17k y de evaluar explícitamente el sesgo en los resultados.

\subsection{Consecuencias y Prevención}

Las tres dimensiones descritas convergen en una conclusión: existe una necesidad concreta de sistemas de apoyo diagnóstico que sean precisos, interpretables y equitativos. La implementación de sistemas de IA para asistencia diagnóstica en dermatología tiene el potencial de impactar positivamente en múltiples niveles. En términos clínicos, puede reducir la variabilidad interobservador y acortar los tiempos de diagnóstico. En términos de accesibilidad, puede democratizar el acceso a evaluaciones de calidad especializada en regiones con déficit de dermatólogos. 
Tschandl et al. (2020) demostraron que la colaboración humano--IA mejora la precisión diagnóstica respecto a cada agente operando de forma aislada \cite{tschandl_human-computer_2020}, mientras que Liu et al. (2024) verificaron que la asistencia de un sistema de aprendizaje profundo incrementó la precisión diagnóstica en más del 33\% tanto para especialistas como para médicos generalistas \cite{liu_deep_learning_aided_2024}. Estos resultados motivan directamente el enfoque del presente trabajo: un VLM capaz de clasificar y explicar sus decisiones diagnósticas.

% ============================================================================
% III. MARCO TEÓRICO TECNOLÓGICO
% ============================================================================
\section{Marco Teórico Tecnológico}

Esta sección sintetiza los fundamentos técnicos que sustentan la propuesta. Para un desarrollo completo de cada concepto, se remite al documento teórico complementario \emph{Documento LLMs, VLLMs}.

\subsection{Definiciones Base: ¿Qué es un Modelo de Visión--Lenguaje?}

Un Modelo de Visión--Lenguaje (VLM, Vision-Language Model) —también referido como Gran Modelo Multimodal (LMM, Large Multimodal Model)— es una arquitectura de red neuronal capaz de procesar simultáneamente información visual (imágenes) y textual (lenguaje natural) para realizar tareas como clasificación, respuesta a preguntas visuales (VQA, Visual Question Answering) y generación de texto condicionado por imágenes \cite{hartsock_vision-language_2024}.

Formalmente, un VLM integra tres componentes principales:
\begin{enumerate}
    \item \textbf{Codificador visual} ($f_v$): extrae representaciones vectoriales de la imagen de entrada. Típicamente un Vision Transformer (ViT, Vision Transformer) \cite{dosovitskiy_image_2021} o un modelo CLIP (Contrastive Language-Image Pretraining) \cite{radford_learning_nodate}.
    \item \textbf{Capa de proyección} ($g$): transforma las representaciones visuales al espacio de incrustaciones del modelo de lenguaje. Generalmente implementada como un Perceptrón Multicapa (MLP, Multilayer Perceptron).
    \item \textbf{Modelo de lenguaje} ($f_l$): un Gran Modelo de Lenguaje (LLM, Large Language Model) decodificador que genera texto autorregesivamente condicionado tanto por las representaciones visuales proyectadas como por el texto de entrada. Ejemplos: LLaMA \cite{touvron_llama_2023}, Vicuna.
\end{enumerate}

La operación del VLM puede expresarse como:
\begin{equation}
    y = f_l\big(g(f_v(x_{\text{img}})),\; x_{\text{text}}\big)
\end{equation}
donde $x_{\text{img}}$ es la imagen de entrada, $x_{\text{text}}$ es la instrucción textual y $y$ es la secuencia de texto generada.

% IMAGEN SUGERIDA: Diagrama general de la arquitectura de un VLM mostrando los tres componentes (codificador visual, proyección, LLM).
% Fuente: Hartsock \& Rasool (2024), Fig. 1 — "Vision-language models for medical report generation and visual question answering: a review", Frontiers in AI.
% DOI: 10.3389/frai.2024.1430984

\subsection{Clasificación del Problema}

A partir de esta definición, las tareas abordadas en esta tesina se clasifican de la siguiente manera:

\begin{itemize}
    \item \textbf{Fase 1 — Clasificación de imágenes:} Problema de clasificación supervisada multiclase. Dado una imagen dermatoscópica $x$, predecir la categoría diagnóstica $\hat{y} \in \{$MEL, NV, BCC, AKIEC, BKL, DF, VASC$\}$.
    \item \textbf{Fase 2 — Generación de reportes:} Problema de generación de texto condicionado por imagen (Image-conditioned Text Generation). Dada una imagen $x_{\text{img}}$ y una instrucción $x_{\text{text}}$, generar un reporte clínico $y = (y_1, y_2, \ldots, y_T)$ de longitud variable.
\end{itemize}

Ambas tareas requieren un modelo capaz de integrar información visual y textual, lo cual motiva la elección de una arquitectura VLM frente a modelos puramente visuales. A continuación se describe el proceso de desarrollo de dicho modelo.

\subsection{Ciclo de Vida de Desarrollo del Modelo}

El desarrollo de un VLM especializado en dermatología sigue un ciclo de vida en etapas progresivas, cada una construyendo sobre los resultados de la anterior:

\begin{enumerate}
    \item \textbf{Pre-entrenamiento del modelo base:} LLaVA-Med fue pre-entrenado por Li et al. (2023) sobre PMC-15M, un corpus de 15 millones de pares imagen--texto biomédicos extraídos de PubMed Central, seguido de 600k pares para alineamiento visual-textual y 60k pares para ajuste de instrucciones \cite{li_llava-med_2023}. Esta etapa no se replica en este trabajo; se utiliza el modelo pre-entrenado disponible públicamente.
    
    \item \textbf{Alineamiento al dominio (Etapa 1):} Se congela el codificador visual y el LLM, entrenando únicamente la capa de proyección MLP con datos dermatológicos (SkinCAP, split de entrenamiento $\sim$2,800 pares imagen-caption EN revisados clínicamente). Esto permite al modelo aprender a mapear conceptos visuales dermatológicos al espacio semántico del LLM sin riesgo de olvido catastrófico (Catastrophic Forgetting).

    \item \textbf{Ajuste de instrucciones (Etapa 2):} Se aplica LoRA sobre las capas de atención del LLM y se entrena conjuntamente con la capa de proyección usando narrativas de razonamiento diagnóstico (SkinCoT EN, split de entrenamiento $\sim$2,130 pares, formato instrucción). Esta es la etapa donde el modelo aprende a generar reportes estructurados con cadena de pensamiento (CoT) clínico.
    
    \item \textbf{Evaluación y validación:} Evaluación cuantitativa con métricas automáticas y análisis cualitativo de reportes generados.
\end{enumerate}

Este paradigma de dos etapas de entrenamiento ha sido validado en múltiples trabajos: LLaVA original \cite{liu_visual_2023}, LLaVA-Med \cite{li_llava-med_2023}, PeFoMed \cite{he_pefomed_2024} y SkinVL \cite{zhao_mm-skin_2025}.

\subsection{Arquitectura Técnica}

La arquitectura concreta que implementa este ciclo de vida se basa en LLaVA-Med v1.5 (7B), cuyos componentes se ilustran en la Figura~\ref{fig:arquitectura} y se detallan en la Tabla~\ref{tab:componentes}.

\begin{figure}[H]
\centering
\begin{tikzpicture}[
    node distance=1.2cm,
    block/.style={rectangle, draw, rounded corners, minimum height=0.8cm, minimum width=2.5cm, align=center, font=\scriptsize},
    arrow/.style={-{Stealth[length=3mm]}, thick}
]
    % Nodos
    \node[block, fill=blue!15] (img) {Imagen\\$224 \times 224$ px};
    \node[block, fill=blue!25, right=of img] (clip) {CLIP ViT-L/14\\(Congelado)};
    \node[block, fill=orange!25, right=of clip] (proj) {MLP Proyección\\(Entrenable)};
    \node[block, fill=green!15, below=1.5cm of img] (text) {Instrucción\\textual};
    \node[block, fill=red!20, right=of proj] (llm) {LLaMA-2 7B\\(LoRA $r$=64)};
    \node[block, fill=purple!15, right=of llm] (out) {Reporte\\clínico};
    
    % Flechas
    \draw[arrow] (img) -- (clip);
    \draw[arrow] (clip) -- (proj);
    \draw[arrow] (proj) -- (llm);
    \draw[arrow] (text) -| (llm);
    \draw[arrow] (llm) -- (out);
\end{tikzpicture}
\caption{Arquitectura del VLM propuesto basado en LLaVA-Med v1.5 (7B). Los componentes en azul permanecen congelados; los componentes en naranja y rojo son entrenables.}
\label{fig:arquitectura}
\end{figure}

% IMAGEN SUGERIDA: Diagrama de la arquitectura original de LLaVA-Med para comparación.
% Fuente: Li et al. (2023), Fig. 2 — "LLaVA-Med: Training a Large Language-and-Vision Assistant for Biomedicine in One Day", NeurIPS.
% Repositorio: https://github.com/microsoft/LLaVA-Med

\begin{table}[H]
\centering
\caption{Especificación técnica de los componentes de la arquitectura.}
\label{tab:componentes}
\begin{tabular}{@{}llll@{}}
\toprule
\textbf{Componente} & \textbf{Modelo} & \textbf{Parámetros} & \textbf{Estado} \\
\midrule
Codificador visual & CLIP ViT-L/14 & $\sim$304M & Congelado \\
Capa de proyección & MLP (2 capas) & $\sim$20M & Entrenable \\
Modelo de lenguaje & LLaMA-2 7B & $\sim$6.7B & LoRA ($\sim$56M entrenables) \\
\midrule
\textbf{Total entrenable} & & $\sim$76M & 1.1\% del total \\
\textbf{VRAM estimada} & & $\sim$14 GB & Compatible con GPU comerciales\\
\bottomrule
\end{tabular}
\end{table}

% ============================================================================
% IV. PROPUESTA DE SOLUCIÓN
% ============================================================================
\section{Propuesta de Solución}

Una vez definidos el problema y la arquitectura general de un VLM, esta sección justifica las decisiones concretas de diseño: qué modelo utilizar, con qué datos entrenarlo, qué técnicas aplicar y cómo evaluar los resultados.

\subsection{Justificación del Modelo}

La selección de LLaVA-Med v1.5 (7B) como modelo base se fundamenta en un análisis de viabilidad de los VLMs del estado del arte según las restricciones de hardware (Usamos como referencia 1 para estimar 1$\times$ NVIDIA RTX 4090, 24 GB VRAM). La Tabla~\ref{tab:viabilidad} presenta este análisis; el modelo seleccionado se resalta en la última fila.

\begin{table}[H]
\centering
\begin{threeparttable}
\caption{Análisis de viabilidad de VLMs dermatológicos en hardware comercial.}
\label{tab:viabilidad}
\small
\begin{tabular}{@{}llcll@{}}
\toprule
\textbf{Modelo} & \textbf{Params} & \textbf{GPU (paper)} & \textbf{¿RTX 4090?} & \textbf{Ref.} \\
\midrule
PanDerm & $\sim$300M enc. & 4$\times$H100 & $\times$ & \cite{yan_panderm_2024} \\
SkinGPT-4 & 13B & 8$\times$A100 & $\triangle$ & \cite{zhou_skingpt4_2024} \\
SkinVL & 7B & 8$\times$3090 & $\checkmark$ (LoRA) & \cite{zhao_mm-skin_2025} \\
MpoxVLM & 7B & 2$\times$RTX 4090 & $\checkmark$ & \cite{xu_mpoxvlm_2024} \\
PeFoMed & 7B & 4$\times$A40 & $\checkmark$ & \cite{he_pefomed_2024} \\
Skin-R1 & 7B & 2$\times$A100 & $\checkmark$ (LoRA) & \cite{kang_skinr1_2025} \\
DermatoLlama & 11B & 4$\times$RTX 4090 & $\triangle$ & \cite{yilmaz_dermatollama_2025} \\
\midrule
\rowcolor{yellow!20}
\textbf{LLaVA-Med v1.5} & \textbf{7B} & \textbf{8$\times$A100\textsuperscript{*}} & \textbf{$\checkmark$ (LoRA)} & \textbf{\cite{li_llava-med_2023}} \\
\bottomrule
\end{tabular}
\begin{tablenotes}
\small
\item $\checkmark$ = viable con LoRA en 1$\times$RTX 4090, $\triangle$ = requiere cuantización agresiva, $\times$ = inviable.
\item \textsuperscript{*} Hardware del paper original; el ajuste fino con LoRA reduce los requisitos.
\end{tablenotes}
\end{threeparttable}
\end{table}

Las razones específicas para la elección de LLaVA-Med 7B son:
\begin{enumerate}
    \item \textbf{Eficiencia en recursos:} Con LoRA ($r$=64), solo $\sim$56M de parámetros son entrenables, requiriendo $\sim$14 GB de VRAM. PeFoMed demostró que esta configuración alcanza resultados competitivos \cite{he_pefomed_2024}.
    \item \textbf{Validado en dermatología:} SkinVL demostró que LLaVA-Med ajustado con datos dermatológicos supera a modelos generalistas (GPT-4V, Gemini) en VQA dermatológico \cite{zhao_mm-skin_2025}.
    \item \textbf{Ecosistema maduro:} Código abierto, ampliamente documentado (\texttt{microsoft/LLaVA-Med}), con múltiples trabajos derivados (SkinVL, MpoxVLM, PeFoMed).
    \item \textbf{Pipeline de dos etapas validado:} El paradigma alineamiento + ajuste de instrucciones ha sido replicado exitosamente en al menos cuatro trabajos médicos \cite{li_llava-med_2023, he_pefomed_2024, zhao_mm-skin_2025, xu_mpoxvlm_2024}.
    \item \textbf{Capacidad generativa:} A diferencia de modelos solo-codificador como CLIP o MONET \cite{kim_monet_2024}, LLaVA-Med puede generar texto libre, requisito indispensable para la generación de reportes.
\end{enumerate}

\subsection{Justificación de los Datasets}

El entrenamiento y evaluación del modelo requiere múltiples tipos de datos: pares imagen--texto para el ajuste fino del VLM, un benchmark de clasificación para evaluación comparativa, un dataset de VQA dermatológico para evaluar la generación de reportes, y un dataset de población local para analizar sesgos. La Tabla~\ref{tab:datasets_comp} presenta una comparación de los datasets candidatos considerados, seguida de la justificación de cada selección.

\begin{table}[H]
\centering
\caption{Comparación de datasets candidatos para el proyecto.}
\label{tab:datasets_comp}
\small
\begin{tabular}{@{}lccccl@{}}
\toprule
\textbf{Dataset} & \textbf{Tamaño} & \textbf{Texto clínico} & \textbf{Diag. confirmado} & \textbf{Diversidad (FST)} & \textbf{Ref.} \\
\midrule
SkinCaRe & 7,041 & Sí (CoT) & Parcial & Alta (Fitz17k+DDI) & \cite{xu_skincare_2025} \\
MM-Skin & 10k+27k VQA & Sí (libros) & No & Media & \cite{zhao_mm-skin_2025} \\
HAM10000 & 10,015 & No & Sí ($>$50\% biopsia) & Baja & \cite{tschandl_ham10000_2018} \\
ISIC 2019 & 25,331 & No & Sí & Baja & \cite{tschandl_ham10000_2018} \\
Derm1M & 1,029,761 & Sí (auto) & No & Variable & \cite{yan_derm1m_2025} \\
HIBA & 1,616 & No & Sí (58\% biopsia) & Media (Lat.Am.) & \cite{piashiba_hiba_2023} \\
\bottomrule
\end{tabular}
\begin{flushleft}
\small FST = Fitzpatrick Skin Type. CoT = Chain-of-Thought. Los datasets seleccionados son SkinCaRe (ajuste fino — SkinCAP en Etapa 1, SkinCoT EN en Etapa 2), MM-Skin (evaluación con referencia + VQA), HAM10000 (clasificación + evaluación sin referencia) y HIBA (evaluación de sesgo).
\end{flushleft}
\end{table}

\subsubsection{Dataset con pares imagen--texto para ajuste fino: SkinCaRe}
\label{sec:skincare_training}

SkinCaRe \cite{xu_skincare_2025} es el dataset principal de ajuste fino. Contiene 7,041 casos dermatológicos divididos en dos subsets complementarios, disponibles como dataset gated en HuggingFace (\texttt{yuhos16/SkinCaRe}):

\begin{itemize}
    \item \textbf{Etapa 1 — Alineamiento (SkinCAP):} $\sim$4,000 pares imagen--caption revisados por dermatólogos en seis ejes clínicos (precisión, seguridad, fundamentación médica, cobertura clínica, coherencia de razonamiento, precisión descriptiva). El split de entrenamiento ($\sim$2,800 pares) se usa para entrenar la capa de proyección MLP. Se utiliza el campo \texttt{caption\_en} del archivo \texttt{skincap\_v240715.xlsx}.
    \item \textbf{Etapa 2 — Instrucción (SkinCoT EN):} 3,041 narrativas de razonamiento diagnóstico jerárquico (CoT), validadas clínicamente. El split de entrenamiento ($\sim$2,130 pares) se usa para el ajuste de instrucciones vía LoRA. Los textos están organizados en \texttt{SkinCoT/EN/<disease\_class>/}.
\end{itemize}

La principal ventaja de SkinCaRe respecto a alternativas como MM-Skin es la presencia de razonamiento CoT explícito en SkinCoT: cada narrativa incluye una cadena de justificación diagnóstica paso a paso, lo cual es precisamente lo que el VLM debe aprender a generar en la Etapa 2. Las anotaciones tienen validación clínica directa, a diferencia de las descripciones de libros de texto de MM-Skin.

\subsubsection{Dataset de clasificación: HAM10000}

HAM10000 \cite{tschandl_ham10000_2018} es un conjunto de 10,015 imágenes dermatoscópicas con diagnóstico confirmado (histopatología en $>$50\% de los casos), organizado en 7 categorías. Es el benchmark estándar de facto para clasificación de lesiones cutáneas, utilizado en múltiples trabajos del estado del arte revisado \cite{aboulmira_skinhealthmate_2024, kang_skinr1_2025, khan_ensemble_2025}.

HAM10000 cumple un doble propósito:
\begin{enumerate}
    \item \textbf{Fase 1:} Benchmark para comparación de modelos de clasificación (CNN vs. ViT vs. VLM).
    \item \textbf{Fase 2:} Al no contener reportes textuales, es el candidato ideal para evaluar la capacidad del VLM de generar reportes sobre imágenes sin texto de referencia. Los metadatos disponibles (diagnóstico, edad, sexo, localización) permiten verificar la coherencia de los reportes generados.
\end{enumerate}

\subsubsection{Dataset de evaluación regional: HIBA}

El dataset del Hospital Italiano de Buenos Aires (HIBA) \cite{piashiba_hiba_2023} contiene 1,616 imágenes (1,270 dermoscópicas y 346 clínicas) recolectadas en Argentina, constituyendo la primera base de datos pública de imágenes dermatológicas de población latinoamericana. El dataset se encuentra disponible públicamente en \url{https://github.com/piashiba/HIBASkinLesionsDataset}. Su incorporación permitirá evaluar el comportamiento del modelo en una población subrepresentada en los benchmarks internacionales, analizando potenciales sesgos asociados a los fototipos de piel predominantes en la región.

\subsubsection{Dataset de evaluación externa: MM-Skin}
\label{sec:mmskin}

MM-Skin \cite{zhao_mm-skin_2025} es el primer dataset multimodal a gran escala específico para dermatología con descripciones textuales profesionales de largo formato. Contiene aproximadamente 10,000 pares imagen--texto extraídos de 15 libros de texto de dermatología, cubriendo tres modalidades de imagen (clínica, dermoscópica y patológica), complementados con más de 27,000 muestras de VQA. En este proyecto, MM-Skin cumple exclusivamente roles de \textbf{evaluación}:

\begin{enumerate}
    \item \textbf{Benchmark de reportes con referencia (Nivel 1):} El split de test proporciona pares imagen--texto de referencia para evaluar la calidad de los reportes generados con métricas automáticas (BLEU-4, ROUGE-L, BERTScore, METEOR, CIDEr). Al ser un dataset externo al entrenamiento, constituye una evaluación más honesta de la capacidad de generalización del modelo.
    \item \textbf{Evaluación de VQA:} Permite evaluar la capacidad del VLM para responder preguntas clínicas sobre imágenes dermatológicas, una tarea complementaria a la generación libre de reportes.
    \item \textbf{Diversidad de modalidad:} A diferencia de HAM10000 (predominantemente dermoscópico), MM-Skin incluye imágenes patológicas y clínicas, lo cual permite evaluar la robustez del modelo ante distintas modalidades.
\end{enumerate}

SkinVL, el modelo entrenado por los autores de MM-Skin sobre LLaVA-Med-v1.5 con los mismos datos de MM-Skin, constituye un punto de referencia relevante para contextualizar los resultados de evaluación.

\subsection{Modelos de Comparación (Baselines)}

El diseño experimental requiere baselines para ambas fases. En la Fase 1 (clasificación) se comparan tres paradigmas arquitectónicos; en la Fase 2 (generación de reportes) se comparan dos VLMs de diferente escala. La Tabla~\ref{tab:baselines} resume todos los modelos de comparación.

\begin{table}[H]
\centering
\caption{Modelos de comparación para las Fases 1 y 2.}
\label{tab:baselines}
\small
\begin{tabular}{@{}lllp{6.5cm}@{}}
\toprule
\textbf{Fase} & \textbf{Rol} & \textbf{Modelo} & \textbf{Justificación} \\
\midrule
1 & CNN baseline & DenseNet121 & Mejor balance exactitud/eficiencia para clasificación de melanoma. Khan et al. reportan 92.30\% de exactitud y AUC de 0.951 \cite{khan_ensemble_2025}. \\
1 & ViT baseline & ViT-B/16 & Misma familia que el codificador visual del VLM (CLIP ViT-L/14). Aísla el efecto del componente visual vs. la integración multimodal. \\
1, 2 & VLM principal & LLaVA-Med 7B & Modelo central de la tesina. 7B parámetros, CLIP ViT-L/14 + LLaMA-2 7B. Referencia para clasificación (vía prompt) y generación de reportes. \\
1, 2 & Mini-VLM & CLIP ViT-B/16 + Phi-3-mini 3.8B & VLM de escala reducida ($\sim$4B params totales). Misma arquitectura tipo LLaVA pero con componentes más pequeños. Evalúa el efecto de la escala del modelo. \\
\bottomrule
\end{tabular}
\end{table}

\subsubsection{Justificación del Mini-VLM}

Para evaluar el impacto de la escala del modelo en el rendimiento, se propone construir un VLM de tamaño reducido siguiendo la misma arquitectura tipo LLaVA pero con componentes más pequeños. La Tabla~\ref{tab:mini_vlm} compara los componentes de ambos modelos.

\begin{table}[H]
\centering
\caption{Comparación de componentes: LLaVA-Med vs. Mini-VLM.}
\label{tab:mini_vlm}
\small
\begin{tabular}{@{}llll@{}}
\toprule
\textbf{Componente} & \textbf{LLaVA-Med 7B} & \textbf{Mini-VLM} & \textbf{Reducción} \\
\midrule
Codificador visual & CLIP ViT-L/14 ($\sim$304M) & CLIP ViT-B/16 ($\sim$86M) & 3.5$\times$ \\
Capa de proyección & MLP (2 capas, $\sim$20M) & MLP (2 capas, $\sim$10M) & 2$\times$ \\
Modelo de lenguaje & LLaMA-2 7B & Phi-3-mini 3.8B & 1.8$\times$ \\
\midrule
\textbf{Total} & $\sim$7.3B & $\sim$3.9B & 1.9$\times$ \\
\textbf{Entrenable (LoRA)} & $\sim$76M & $\sim$42M & 1.8$\times$ \\
\textbf{VRAM estimada} & $\sim$14 GB & $\sim$8 GB & 1.75$\times$ \\
\bottomrule
\end{tabular}
\end{table}

La elección de Phi-3-mini (3.8B parámetros) \cite{abdin_phi3_2024} como LLM del Mini-VLM se fundamenta en tres razones: (1) es un modelo sub-4B con un gran rendimiento en benchmarks de razonamiento y generación de texto a la fecha de su publicación; (2) su arquitectura tipo decoder-only es compatible con la integración de tokens visuales mediante proyección MLP, siguiendo el mismo paradigma que LLaVA; y (3) al ser desarrollado por Microsoft (al igual que LLaVA-Med), se espera buena compatibilidad técnica del ecosistema. La combinación con CLIP ViT-B/16 (en lugar de ViT-L/14) reduce la dimensionalidad de las representaciones visuales de 1024 a 768, acorde con la capacidad del LLM más pequeño.

El Mini-VLM se entrena con el mismo pipeline de dos etapas, los mismos datos (MM-Skin) y los mismos hiperparámetros de LoRA que LLaVA-Med, variando únicamente la escala de los componentes. Esta configuración de variables controladas permite atribuir las diferencias de rendimiento exclusivamente al factor de escala. La comparación es relevante porque un modelo significativamente más pequeño que mantenga rendimiento aceptable sería preferible para despliegue en entornos clínicos con recursos computacionales limitados. 

\subsection{Técnicas y Algoritmos}

\subsubsection{Ajuste fino con LoRA (Low-Rank Adaptation)}

LoRA \cite{hu_lora_2021} es una técnica de Ajuste Fino Eficiente en Parámetros (PEFT, Parameter-Efficient Fine-Tuning) que inyecta matrices de bajo rango en las capas de atención del LLM. Para una matriz de pesos $W_0 \in \mathbb{R}^{d \times k}$, LoRA introduce la descomposición:
\begin{equation}
    W = W_0 + \Delta W = W_0 + BA, \quad B \in \mathbb{R}^{d \times r}, \; A \in \mathbb{R}^{r \times k}, \; r \ll \min(d,k)
\end{equation}
donde solo $A$ y $B$ son entrenables, reduciendo los parámetros de $d \times k$ a $r \times (d + k)$.

\textbf{Configuración propuesta:} rango $r = 64$, $\alpha = 128$, dropout = 0.05, aplicado a las matrices $W_q$ y $W_v$ de las capas de atención del LLM. Esta configuración es consistente con la utilizada por PeFoMed \cite{he_pefomed_2024}.

\textbf{QLoRA:} Si la VRAM resulta insuficiente, se aplicará cuantización de 4 bits (QLoRA) \cite{dettmers_qlora_2023} sobre el LLM base, reduciendo el consumo.

Seis de los doce VLMs dermatológicos revisados en el estado del arte emplean LoRA como técnica principal de ajuste fino: SkinVL \cite{zhao_mm-skin_2025}, MpoxVLM \cite{xu_mpoxvlm_2024}, PeFoMed \cite{he_pefomed_2024}, Chakrabarti et al. \cite{chakrabarti_automated_2025}, Skin-R1 \cite{kang_skinr1_2025} y DermatoLlama \cite{yilmaz_dermatollama_2025}.

\subsubsection{Entrenamiento en dos etapas}

El pipeline de dos etapas se aplica de manera idéntica a LLaVA-Med y al Mini-VLM, variando únicamente los componentes internos (codificador visual y LLM). Esta simetría permite una comparación justa del efecto de la escala.

\textbf{Etapa 1 — Alineamiento visual-textual:}
\begin{itemize}
    \item Se congela el codificador visual (CLIP ViT-L/14) y el LLM.
    \item Se entrena la capa de proyección MLP con SkinCAP captions EN, split de entrenamiento ($\sim$2,800 pares, validados clínicamente).
    \item Objetivo: aprender el mapeo de características visuales dermatológicas al espacio semántico del LLM.
\end{itemize}

\textbf{Etapa 2 — Ajuste de instrucciones:}
\begin{itemize}
    \item Se congela el codificador visual.
    \item Se entrena el LLM vía LoRA + la capa de proyección con SkinCoT EN, split de entrenamiento ($\sim$2,130 pares, formato instrucción con CoT diagnóstico).
    \item Formato de instrucción: \texttt{``<image>\textbackslash n Describe esta lesión cutánea con un reporte clínico estructurado que incluya: hallazgos, impresión diagnóstica y diagnóstico diferencial.''}
\end{itemize}

\subsubsection{Preprocesamiento de datos}

\begin{table}[H]
\centering
\caption{Pipeline de preprocesamiento de datos.}
\label{tab:preproc}
\begin{tabular}{@{}llp{7cm}@{}}
\toprule
\textbf{Paso} & \textbf{Detalle} & \textbf{Justificación} \\
\midrule
Redimensionamiento & 224$\times$224 px & Resolución nativa de CLIP ViT-L/14 \\
Normalización & Media y std de ImageNet & Requerido por codificador pre-entrenado \\
Aumento (Fase 1) & Rotación $\pm$30°, flip, color jitter & Mitigación de desbalance de clases \cite{daneshjou_disparities_2022, aboulmira_skinhealthmate_2024} \\
Splits & 70\%/15\%/15\% & Estratificado por diagnóstico \\
\bottomrule
\end{tabular}
\end{table}

\subsection{Métricas de Evaluación}

Las métricas se dividen según la fase experimental. Para cada métrica se indica en qué trabajos del estado del arte ha sido utilizada, lo cual justifica su selección.

\subsubsection{Fase 1: Clasificación Diagnóstica}

\begin{table}[H]
\centering
\caption{Métricas de evaluación para clasificación diagnóstica.}
\label{tab:metricas_clas}
\begin{tabular}{@{}lp{5.5cm}p{5.5cm}@{}}
\toprule
\textbf{Métrica} & \textbf{Descripción} & \textbf{Utilizada en} \\
\midrule
Exactitud (Accuracy) & Proporción de predicciones correctas & \cite{yan_panderm_2024, yilmaz_dermatollama_2025, xu_mpoxvlm_2024, groh_fitzpatrick17k_2021} \\
AUC-ROC & Capacidad discriminativa independiente del umbral & \cite{yan_panderm_2024, yan_derm1m_2025, kim_monet_2024, daneshjou_disparities_2022} \\
F1-Score (Ponderado) & Media armónica ponderada; maneja desbalance de clases & \cite{yan_panderm_2024, xu_mpoxvlm_2024, khan_ensemble_2025} \\
Sensibilidad (Recall) & Proporción de positivos correctamente identificados & \cite{khan_ensemble_2025, daneshjou_disparities_2022} \\
Especificidad & Proporción de negativos correctamente identificados & \cite{kim_monet_2024, daneshjou_disparities_2022} \\
\bottomrule
\end{tabular}
\end{table}

\subsubsection{Fase 2: Generación de Reportes}

\begin{table}[H]
\centering
\caption{Métricas de evaluación para generación de reportes clínicos.}
\label{tab:metricas_gen}
\begin{tabular}{@{}llp{4.5cm}p{4cm}@{}}
\toprule
\textbf{Métrica} & \textbf{Tipo} & \textbf{Descripción} & \textbf{Utilizada en} \\
\midrule
BLEU-4 & Léxica & Precisión de 4-gramas & \cite{yilmaz_dermatollama_2025, zhao_mm-skin_2025} \\
ROUGE-L & Léxica & Subsecuencia común más larga & \cite{yilmaz_dermatollama_2025, chakrabarti_automated_2025, zhao_mm-skin_2025, he_pefomed_2024} \\
BERTScore F1 & Semántica & Similitud coseno entre incrustaciones contextuales & \cite{chakrabarti_automated_2025, zhao_mm-skin_2025} \\
METEOR & Léxica + Sem. & Precisión, recall y penalización por fragmentación & \cite{zhao_mm-skin_2025, he_pefomed_2024} \\
CIDEr & Léxica (TF-IDF) & Similitud ponderada por relevancia del n-grama & \cite{he_pefomed_2024} \\
\bottomrule
\end{tabular}
\end{table}

Adicionalmente, se incorporará una evaluación de \textbf{Coincidencia de Entidades Clínicas} (Clinical Entity Matching) basada en Reconocimiento de Entidades Nombradas (NER, Named Entity Recognition), que verifica si los reportes generados mencionan conceptos clínicos clave como asimetría, bordes irregulares, variaciones cromáticas y diámetro. Esta métrica, inspirada en la evaluación de características del método de los 7 puntos utilizada por Chakrabarti et al. \cite{chakrabarti_automated_2025}, aporta una dimensión de evaluación clínicamente relevante que las métricas puramente léxicas no capturan.

\subsubsection{Evaluación Comparativa de Reportes}

Mientras que la Fase 1 compara modelos exclusivamente en clasificación, la Fase 2 requiere un protocolo específico para evaluar y comparar la generación de reportes entre LLaVA-Med y el Mini-VLM. Este se estructura en tres niveles:

\textbf{Nivel 1 — Evaluación automática con referencia (MM-Skin test split).} Ambos VLMs generan reportes sobre el split de test de MM-Skin ($\sim$20\% del dataset). Los reportes generados se comparan contra los textos de referencia (ground truth) utilizando las métricas léxicas y semánticas de la Tabla~\ref{tab:metricas_gen}: BLEU-4, ROUGE-L, BERTScore, METEOR y CIDEr. MM-Skin actúa como benchmark externo: el modelo fue entrenado con SkinCaRe y es evaluado sobre MM-Skin, lo que constituye una evaluación de generalización más rigurosa que evaluar sobre datos de entrenamiento.

\textbf{Nivel 2 — Evaluación automática sin referencia (HAM10000 + HIBA).} Como HAM10000 y HIBA carecen de reportes textuales, los reportes generados se evalúan mediante: (a) Coincidencia de Entidades Clínicas, verificando que el reporte mencione hallazgos coherentes con el diagnóstico conocido del paciente (metadatos disponibles: diagnóstico, edad, sexo, localización); y (b) consistencia diagnóstica, comprobando si la impresión diagnóstica del reporte coincide con la etiqueta de clasificación.

\textbf{Nivel 3 — Análisis cualitativo.} Se selecciona una muestra estratificada de reportes (al menos 5 por categoría diagnóstica y por modelo) para un análisis cualitativo que evalúa: coherencia clínica, completitud del reporte (hallazgos, impresión, diagnóstico diferencial), presencia de alucinaciones y legibilidad. Este nivel, aunque no cuantitativo, es esencial para detectar modos de fallo que las métricas automáticas no capturan, como se evidenció en la evaluación humana reportada por Chakrabarti et al. \cite{chakrabarti_automated_2025}.

La Tabla~\ref{tab:eval_reportes} resume este protocolo.

\begin{table}[H]
\centering
\caption{Evaluación comparativa de reportes (Fase 2).}
\label{tab:eval_reportes}
\small
\begin{tabular}{@{}llll@{}}
\toprule
\textbf{Nivel} & \textbf{Dataset} & \textbf{Métricas} & \textbf{Modelos} \\
\midrule
1 (con referencia) & MM-Skin (test, $\sim$2.2k) & BLEU-4, ROUGE-L, BERTScore, METEOR, CIDEr & LLaVA-Med, Mini-VLM \\
2 (sin referencia) & HAM10000 & Entidades clínicas, consistencia diagnóstica & LLaVA-Med, Mini-VLM \\
2 (sin referencia) & HIBA & Entidades clínicas, consistencia diagnóstica & LLaVA-Med, Mini-VLM \\
3 (cualitativo) & Muestra cruzada & Coherencia, completitud, alucinaciones & LLaVA-Med, Mini-VLM \\
\bottomrule
\end{tabular}
\end{table}

% ============================================================================
% V. ESTADO DEL ARTE / TRABAJOS RELACIONADOS
% ============================================================================
\section{Estado del Arte}

\subsection{Evolución de la IA en Dermatología}

La aplicación de inteligencia artificial al diagnóstico dermatológico ha experimentado una evolución acelerada que puede organizarse en tres fases. Esta progresión refleja tanto los avances en las arquitecturas de aprendizaje profundo como la creciente disponibilidad de datos dermatológicos, y permite situar la contribución del presente trabajo en la frontera actual de investigación.

\textbf{Fase I — Modelos puramente visuales (2017--2022).} El trabajo seminal de Esteva et al. (2017) demostró por primera vez que una CNN podía alcanzar rendimiento comparable al de dermatólogos certificados \cite{esteva_dermatologist_2017}. A partir de este hito, la investigación fue dominada por CNNs pre-entrenadas en ImageNet y ajustadas para clasificación dermatológica. Arquitecturas como DenseNet121, EfficientNet y ResNet se consolidaron como baselines, con resultados que oscilan entre 85\% y 97\% de exactitud dependiendo del dataset y la complejidad de la tarea \cite{khan_ensemble_2025, aboulmira_skinhealthmate_2024}. La principal limitación de estos modelos es su naturaleza de ``caja negra'': producen una etiqueta diagnóstica sin justificación textual.

\textbf{Fase II — Modelos multimodales y fundacionales (2023--2024).} Para superar esta limitación, la emergencia de modelos como CLIP y LLaVA permitió la integración de visión y lenguaje. Trabajos como MONET \cite{kim_monet_2024} y SkinCon \cite{daneshjou_skincon_2022} introdujeron el uso de conceptos semánticos para auditoría e interpretabilidad. LLaVA-Med \cite{li_llava-med_2023} estableció el paradigma de alineamiento + ajuste de instrucciones para VLMs biomédicos, demostrando que un modelo de 7B parámetros podía ser entrenado en menos de 15 horas con resultados competitivos en tareas biomédicas.

\textbf{Fase III — VLMs especializados en dermatología (2024--2025).} Los trabajos más recientes se enfocan en adaptar VLMs específicamente para dermatología: SkinVL \cite{zhao_mm-skin_2025} con fine-tuning de LLaVA-Med, SkinGPT-4 \cite{zhou_skingpt4_2024} como sistema interactivo, PeFoMed \cite{he_pefomed_2024} con fine-tuning eficiente, y Skin-R1 \cite{kang_skinr1_2025} con razonamiento por refuerzo. Es en esta fase donde se inserta el presente trabajo, con la particularidad de combinar clasificación comparativa y generación de reportes en un pipeline unificado.

\subsection{Benchmarking: Trabajos de Referencia}

La Tabla~\ref{tab:sota} resume los VLMs dermatológicos más relevantes del estado del arte. Para cada modelo se reporta su mejor resultado publicado junto con la métrica y la técnica de ajuste fino utilizada. Dado que los modelos difieren en tareas y datasets de evaluación, los valores numéricos no son directamente comparables entre filas; se incluyen para ilustrar el nivel de rendimiento alcanzado en cada caso.

\begin{longtable}{@{}cllp{2.5cm}lll@{}}
\caption{Resumen comparativo de VLMs dermatológicos del estado del arte. La columna ``Técnica FT'' indica la estrategia de ajuste fino; ``2 etapas'' se refiere al paradigma de alineamiento seguido de ajuste de instrucciones.} \label{tab:sota} \\
\toprule
\textbf{\#} & \textbf{Modelo} & \textbf{Arquitectura} & \textbf{Técnica FT} & \textbf{Métrica} & \textbf{Mejor resultado} \\
\midrule
\endfirsthead
\toprule
\textbf{\#} & \textbf{Modelo} & \textbf{Arquitectura} & \textbf{Técnica FT} & \textbf{Métrica} & \textbf{Mejor resultado} \\
\midrule
\endhead
1 & PanDerm & ViT-L + CLIP & SSL (MAE+CLIP) & AUROC & 0.964 \\
2 & DermLIP & ViT-B/16 + Text Enc. & Contrastivo & Accuracy & 56.07\% \\
3 & DermatoLlama & Llama 3.2-11B & LoRA & ROUGE-L & 0.93 \\
4 & LLaVA-Med & CLIP + Vicuna 7B & 2 etapas & Accuracy & 84.19\% \\
5 & SkinVL & LLaVA-Med-v1.5-7B & LoRA + 2 etapas & BLEU-4 & 22.04 \\
6 & MpoxVLM & CLIP + LLaMA-2-7B & LoRA + 2 etapas & Accuracy & 90.38\% \\
7 & PeFoMed & EVA-ViT + LLaMA-2 7B & LoRA + 2 etapas & Accuracy & 87.5\% \\
8 & Chakrabarti & CLIP + LLaMA-2 13B & KD + LoRA & BERTScore & 0.9025 \\
9 & SkinGPT-4 & ViT + Llama-2-13B & 2 etapas & Ac. clínico & 80.63\% \\
10 & Skin-R1 & Qwen2.5-VL-7B & SFT + RL (GRPO) & Accuracy & 63.85\% \\
11 & MONET & ViT-L/14 (contrastivo) & Contrastivo & AUROC & 0.881 \\
12 & Khan et al. & Ensemble + LLaMA-3 70B & Votación + Prompt & Accuracy & 97.67\% \\
\bottomrule
\end{longtable}

% IMAGEN SUGERIDA: Gráfico de barras comparando resultados de los modelos de 7B con LoRA (SkinVL, MpoxVLM, PeFoMed, LLaVA-Med) en tareas de clasificación.
% Los datos se extraen directamente de la Tabla anterior.

De este análisis se extraen las siguientes observaciones que fundamentan las decisiones de diseño:
\begin{itemize}
    \item Los modelos de 7B parámetros con LoRA (SkinVL, MpoxVLM, PeFoMed) logran resultados competitivos frente a modelos de mayor tamaño, validando la viabilidad del enfoque propuesto.
    \item El entrenamiento en dos etapas (alineamiento + ajuste de instrucciones) es el paradigma más adoptado por los VLMs generativos del estado del arte: lo utilizan LLaVA-Med, SkinVL, MpoxVLM, PeFoMed y SkinGPT-4, como se refleja en la columna ``Técnica FT'' de la Tabla~\ref{tab:sota}.
    \item BERTScore y ROUGE-L son las métricas más frecuentes para evaluación de generación de reportes.
    \item Ningún trabajo previo combina clasificación comparativa (CNN vs. ViT vs. VLM) con generación de reportes en un pipeline unificado, lo cual constituye la contribución diferencial de esta tesina.
    \item La mayoría de los modelos revisados utilizan LLMs de 7B+ parámetros. Esta tesina incluye una comparación con un VLM de escala reducida ($\sim$4B) para evaluar si la eficiencia del modelo puede mantenerse a menor costo computacional.
\end{itemize}

% ============================================================================
% VI. PLAN DE EJECUCIÓN
% ============================================================================
\section{Plan de Ejecución}

El desarrollo experimental se organiza en cuatro fases con hitos verificables. La duración estimada de cada fase se detalla en el cronograma oficial de la tesina; aquí se describen las actividades y entregables correspondientes.

\begin{table}[H]
\centering
\caption{Fases del desarrollo experimental y sus entregables.}
\label{tab:cronograma}
\small
\begin{tabular}{@{}clp{8.5cm}l@{}}
\toprule
\textbf{Fase} & \textbf{Nombre} & \textbf{Actividades} & \textbf{Entregable} \\
\midrule
\rowcolor{phase0!15}
0 & Preparación & Descarga de datasets (HAM10000, SkinCaRe via HuggingFace, MM-Skin, HIBA), configuración de entorno, análisis exploratorio de datos, preprocesamiento y definición de splits estratificados. & Datos listos \\
\rowcolor{phase1!15}
1 & Clasificación & Entrenar DenseNet121 y ViT-B/16 (transfer learning desde ImageNet) sobre HAM10000. Adaptar LLaVA-Med y Mini-VLM para clasificación vía prompt. Evaluación comparativa con métricas de la Tabla~\ref{tab:metricas_clas} & Tabla comparativa \\
\rowcolor{phase2!15}
2 & Generación & Entrenamiento de dos etapas para LLaVA-Med y Mini-VLM (alineamiento con SkinCAP + instruction tuning con SkinCoT EN). Evaluación de reportes según protocolo de la Tabla~\ref{tab:eval_reportes} sobre MM-Skin (test, benchmark externo), HAM10000 y HIBA & Reportes generados \\
\rowcolor{phase3!15}
3 & Integración & Análisis cruzado clasificación/generación. Estudios de ablación (rango LoRA, escala del modelo LLaVA-Med vs. Mini-VLM, LoRA vs. QLoRA). Documentación y escritura final de la tesina & Tesina completa \\
\bottomrule
\end{tabular}
\end{table}

\subsection{Riesgos y Mitigaciones}

\begin{table}[H]
\centering
\caption{Matriz de riesgos y estrategias de mitigación.}
\label{tab:riesgos}
\begin{tabular}{@{}p{4cm}lp{6cm}@{}}
\toprule
\textbf{Riesgo} & \textbf{Prob.} & \textbf{Mitigación} \\
\midrule
VRAM insuficiente con LoRA & Media & Activar QLoRA 4-bit ($\sim$10 GB) \\
SkinCaRe inaccesible & \textbf{Resuelto} & SkinCaRe disponible en HuggingFace (\texttt{yuhos16/SkinCaRe}, dataset gated). Descarga automatizada con \texttt{snapshot\_download} vía \texttt{HF\_TOKEN}. Ver Sección~\ref{sec:skincare_training}. \\
Tiempo de entrenamiento excesivo & Media & Reducir épocas, subsets estratificados, early stopping \\
Alucinaciones en reportes & Alta & Clinical Entity Matching + revisión cualitativa manual \\
\bottomrule
\end{tabular}
\end{table}

\subsection{Infraestructura}

\begin{itemize}
    \item \textbf{Hardware local:} GPU en pc personal NVIDIA RTX 4070 12GB para pruebas locales.
    \item \textbf{Hardware Primario:} Cluster TOKO (ICB-ITIC, UNCuyo) y Mendieta (CCAD, UNC).
    \item \textbf{Repositorios base:} \texttt{microsoft/LLaVA-Med}, \texttt{microsoft/Phi-3-mini-4k-instruct}, \texttt{yuhos16/SkinCaRe} (HuggingFace), \texttt{ZwQ803/MM-Skin}, \texttt{piashiba/HIBASkinLesionsDataset}.
\end{itemize}

% ============================================================================
% BIBLIOGRAFÍA
% ============================================================================
\newpage
\printbibliography[title=Referencias]

\end{document}
